\chapter{Estrutura e organização}

De acordo com a
\href{https://transparencia.cfp.org.br/crp06/wp-content/uploads/sites/29/2022/09/Resolucao-03-22-PECS-f-1-assinada.pdf}{Resolução
  CRP nº~03, de 29 de agosto de 2022}, a estrutura administrativa do
Conselho possui duas instâncias decisórias: gerências e coordenações.

A Coordenação de Comunicação está subordinada à Gerência de Relações
Institucionais. É responsável por elaborar e executar o planejamento de
comunicação definido periodicamente com a Comissão de Comunicação
Institucional (ComCom).

\section{Comissão de Comunicação Institucional
  (ComCom)}

A Comissão de Comunicação Institucional (ComCom) é uma comissão
permanente, que funciona como instância deliberativa e consultiva para
os assuntos relativos à comunicação do CRP~SP.

\subsection{Recursos}

A ComCom é composta de pelo menos duas conselheiras, tendo uma delas
como coordenadora.

\subsection{Atribuições}

De acordo com o
\href{https://correio.crpsp.org.br/service/home/~/?auth=co&loc=pt_BR&id=4201&part=6}{Regimento
  Interno do CRP~SP}, compete à ComCom:

\begin{enumerate}

  \item
        acompanhar e aprovar as ações da Unidade de Comunicação Institucional
        do CRP-06, informando o Plenário;
  \item
        responsabilizar-se pela comunicação institucional interna e externa;
  \item
        responsabilizar-se pela produção e divulgação dos eventos
        institucionais externos de âmbito estadual;
  \item
        dar diretrizes para a produção e divulgação dos eventos do CRP-06;
  \item
        apresentar ao Plenário a identidade visual e o plano de comunicação
        das campanhas, ações e unidades do CRP-06;
  \item
        elaborar a Política de Comunicação Institucional para aprovação do
        Plenário e fazer cumprir suas previsões;
  \item
        acompanhar o adequado preenchimento e publicação do Portal da
        Transparência de acordo com o disposto na Lei de Acesso à Informação
        (LAI; Lei nº~12.527/2011) e na Resolução CFP nº~20/2018, ou outras que
        vierem a substituí-las;
  \item
        propor pautas e acompanhar a elaboração do Jornal Psi para aprovação
        do Plenário;
  \item
        criar o plano anual de comunicação institucional do CRP-06 em
        observância às priorizações do planejamento estratégico;
  \item
        acompanhar e alimentar o site e redes sociais do CRP-06, encaminhando
        propostas e/ou denúncias, que ferem a imagem da autarquia, para
        providências da Diretoria, quando necessário;
  \item
        realizar reuniões semanais com suas/seus membras/os para planejar,
        encaminhar e avaliar as ações;
  \item
        articular com a Comissão de Comunicação do Conselho Federal de
        Psicologia;
  \item
        supervisionar a edição das publicações do CRP-06;
  \item
        coordenar a divulgação das ações do CRP-06;
  \item
        acolher e encaminhar à unidade responsável para que responda às
        solicitações e consultas de pessoas físicas e jurídicas em relação às
        informações de domínio comum da entidade no site e redes sociais, bem
        como aquelas correlatas que lhe sejam atribuídas pelo Plenário;
  \item
        organizar e manter atualizados os canais de comunicação do CRP-06;
  \item
        orientar o diálogo com a imprensa e outros canais externos para
        participação do CRP-06 em suas pautas;
  \item
        garantir a aplicação da linguagem gendrada, antirracista e inclusiva
        no CRP-06, em seus canais de comunicação internos e externos;
  \item
        acompanhar e apresentar ao Plenário os resultados quantitativos e
        qualitativos da Comissão de Comunicação;
  \item
        aprovar custos envolvendo fornecedoras/es e aquisição de outros
        recursos para a plena realização do planejamento e melhoria da área,
        acompanhando as licitações demandadas pela área;
  \item
        propor, avaliar e encaminhar para aprovação do Plenário, sempre que
        necessário, produtos de mídia que possam ser usados para comunicação
        com as profissionais de Psicologia.
\end{enumerate}

\section{Gerência de Relações
  Institucionais}

Tem a finalidade de gerir, avaliar e acompanhar a execução das
atividades das coordenações subordinadas, colaborando com a diretoria e
exercendo as competências que lhe forem delegadas.

\subsection{Recursos}

A Gerência de Relações Institucionais é ocupada por uma trabalhadora em
cargo de comissão.

\subsection{Atribuições}

Compete à Gerência de Relações Institucionais, segundo a
\href{https://transparencia.cfp.org.br/crp06/wp-content/uploads/sites/29/2022/09/Resolucao-03-22-PECS-f-1-assinada.pdf}{Resolução
  CRP nº~03, de 29 de agosto de 2022}:

\begin{enumerate}

  \item
        executar as atividades relacionadas à comunicação e articulação com os
        segmentos da sociedade envolvidos com os serviços fiscalizados pelo
        CRP-06;
  \item
        articular e desenvolver o relacionamento do CRP-06 com órgãos de
        interesse, tais como universidades, organismos nacionais e
        internacionais e organizações da sociedade civil;
  \item
        executar atividades de divulgação de conteúdos, participação e
        realização de eventos relacionados ao trabalho desempenhado pelo
        CRP-06 e de interesse a suas/seus inscritas/os, divulgar conteúdos de
        interesse do CRP-06;
  \item
        mapear, articular e acompanhar possíveis parcerias e convênios com
        órgãos relevantes aos fins do CRP-06, submetendo as propostas à
        Diretoria;
  \item
        atuar de forma articulada com as demais Gerências responsabilizando-se
        pelas deliberações relativas aos processos e à gestão do trabalho no
        cumprimento das metas institucionais;
  \item
        organizar, articular, acompanhar, avaliar e monitorar as atividades
        das unidades sob sua responsabilidade, utilizando-se de indicadores,
        quando necessário, para verificação das necessidades e garantia do bom
        desenvolvimento dos trabalhos; e
  \item
        desempenhar outras atribuições correlatas de nível similar.
\end{enumerate}

\section{Coordenação de
  Comunicação}

A Coordenação de Comunicação é a instância executiva da comunicação do
CRP~SP.

\subsection{Recursos}

A Coordenação de Comunicação é composta de:

\begin{enumerate}

  \item
        uma coordenadora;
  \item
        duas analistas em gestão de comunicação;
  \item
        quatro profissionais de suporte administrativo (PSA);
  \item
        duas estagiárias.
\end{enumerate}

\subsection{Atribuições}

De acordo com a
\href{https://transparencia.cfp.org.br/crp06/wp-content/uploads/sites/29/2022/09/Resolucao-03-22-PECS-f-1-assinada.pdf}{Resolução
  CRP nº~03, de 29 de agosto de 2022}, compete à Coordenação de
Comunicação:

\begin{enumerate}

  \item
        planejar, organizar, coordenar e divulgar as ações referentes à
        integração das estratégias institucionais com a categoria e a
        sociedade;
  \item
        construir, desenvolver e implantar estratégias com foco na aproximação
        com a categoria e a sociedade;
  \item
        garantir a divulgação das atividades e posicionamentos institucionais
        do CRP-06 para a categoria e a sociedade;
  \item
        gerenciar os servidores de \emph{e-mail} e hospedagem de \emph{sites}
        do CRP-06;
  \item
        criar, manter e atualizar as aplicações \emph{web};
  \item
        desenvolver trabalhos de artes gráficas que sejam demandados pelo
        CRP-06, acompanhando as produções e a aprovação do material
        apresentado pelos fornecedores (prova impressa, qualidade de impressão
        e material produzido);
  \item
        fazer a cobertura de eventos com a utilização de filmagens,
        fotografias ou quaisquer outros meios necessários, assim como promover
        a sonorização, projeção de mídia e edição subsequentes necessárias;
  \item
        produzir e diagramar publicações do CRP-06 (impressas ou eletrônicas),
        para divulgação às/aos psicólogas/os, contemplando assuntos de seus
        interesses, quando for o caso;
  \item
        assessorar e desenvolver políticas de relações públicas,
        relacionando-se com as áreas de formação de opinião pública e
        diferentes mídias;
  \item
        desenvolver e executar políticas de assessoria de imprensa,
        acompanhando o alcance e impacto das estratégias de divulgação,
        comunicação e diálogo com a categoria e a sociedade;
  \item
        propor estratégia e divulgar ações técnico-políticas e posicionamentos
        do Conselho visando qualificar e publicizar o CRP-06;
  \item
        divulgar interna e externamente as atividades desenvolvidas pelas
        unidades organizacionais do CRP-06;
  \item
        coordenar e acompanhar a realização de cobertura jornalística
        técnico-política do CRP-06;
  \item
        desenvolver e manter atualizado o banco de imagens e fontes do CRP-06;
  \item
        criar e atualizar \emph{mailing} de contatos de pessoas físicas e
        entidades do CRP-06;
  \item
        realizar todos os tipos de eventos (seminários, congressos, fóruns,
        debates e simpósios, entre outros) do CRP-06, contemplando todas as
        atividades necessárias que percorrem as etapas de planejamento,
        organização, realização e pós-evento;
  \item
        fazer a cobertura de eventos com a utilização de filmagens,
        fotografias ou quaisquer outros meios necessários, assim como promover
        a sonorização, projeção de mídia e edição subsequentes necessárias;
  \item
        manter e organizar os equipamentos audiovisuais do auditório do
        CRP-06;
  \item
        apoiar a realização de eventos indicados pelo CRP-06 para outras
        entidades;
  \item
        desenvolver e manter atualizado o banco de imagens e fontes do CRP-06;
  \item
        realizar exposições do CRP-06 em espaços externos ao CRP-06 em toda a
        sua jurisdição;
  \item
        operacionalizar concursos culturais do CRP-06 no estado de São Paulo;
  \item
        operacionalizar o deslocamento e hospedagem, se necessário, de
        coordenadoras/es, palestrantes e convidadas/os para eventos do CRP-06;
  \item
        pesquisar fornecedores e elaborar cotações para locação de serviços
        e/ou compra de materiais para eventos do CRP-06;
  \item
        pesquisar e apresentar, para avaliação das instâncias superiores,
        espaços para a realização de eventos, analisando as adequações de
        localização e infraestrutura;
  \item
        operar a mesa de som da cabine do auditório do CRP-06, incluindo
        microfones, DVD \emph{players} e demais equipamentos, realizando,
        inclusive, gravação e edição de áudio dos eventos;
  \item
        atender as solicitações diversas das/os psicólogas/os, respondendo por
        \emph{e-mail}, por telefone e pessoalmente, sobre informações de
        eventos do CRP-06;
  \item
        oferecer apoio administrativo na participação do CRP-06 em eventos
        promovidos por outras organizações, fazendo cotações de espaço,
        materiais, equipamentos e participando da promoção de publicações e/ou
        serviços em feiras e outros eventos no estado de São Paulo, prestando
        orientação às/aos visitantes e/ou participantes;
  \item
        criar subsídios para a elaboração de documentos para apoio e
        patrocínio para a realização de eventos realizados pelo CRP-06;
  \item
        identificar e criar estratégias para utilização de novos meios de
        comunicação;
  \item
        desempenhar outras atribuições correlatas de nível similar.
\end{enumerate}

\subsubsection{Atribuições funcionais}

Compete à
\href{https://transparencia.cfp.org.br/crp06/wp-content/uploads/sites/29/2022/09/Resolucao-03-22-PECS-Anexo-II-Regulamentacao-Plano-de-Empregos-Carreiras-e-Salarios-f-assinada.pdf}{\textbf{coordenadora}}
a gestão de processos e equipes de trabalho da Coordenação de
Comunicação do CRP~SP.

Compete às
\href{https://transparencia.cfp.org.br/crp06/wp-content/uploads/sites/29/2022/09/Resolucao-03-22-PECS-Anexo-II-Regulamentacao-Plano-de-Empregos-Carreiras-e-Salarios-f-assinada.pdf}{\textbf{analistas
    em gestão de comunicação}}:

\begin{enumerate}

  \item
        selecionar, revisar, preparar e distribui materiais para publicação;
  \item
        fotografar e gravar imagens jornalísticas;
  \item
        editar publicações impressas e eletrônicas;
  \item
        selecionar, divulgar e arquivar a comunicação feita a respeito do
        Conselho nos meios impressos e eletrônicos;
  \item
        manter contato com a imprensa externa fornecendo dados, materiais e
        marcando entrevistas;
  \item
        implantar ações de relações públicas e assessoria de imprensa;
  \item
        organizar eventos internos e externos;
  \item
        preparar, organizar, coordenar e realizar o cerimonial; e
  \item
        desempenhar outras atribuições correlatas de nível similar.
\end{enumerate}

Compete às
\href{https://transparencia.cfp.org.br/crp06/wp-content/uploads/sites/29/2023/03/Resolucao-CRP-002-2023.pdf}{\textbf{profissionais
    de suporte administrativo}} \textbf{(\emph{design} gráfico)}:

\begin{enumerate}

  \item
        preparar, organizar, elaborar, atualizar e conferir documentos,
        relatórios, gráficos e tabelas bem como prestar atendimento ao
        profissional e ao público em geral na área em que atua;
  \item
        executar, sob supervisão, atividades de preparo técnico de comunicação
        interna e externa, gráficas utilizando a suíte de aplicativos de
        aplicativos Adobe (InDesign, Illustrator, Photoshop), CorelDraw e
        aplicativos similares;
  \item
        preparar composições gráficas, diagramação, desenvolver marcas,
        materiais impressos, papelaria, sinalização, interface de web,
        infográficos e ilustrações;
  \item
        definir especificações técnicas para produção de materiais impressos;
  \item
        acompanhar a produção dos materiais analisando e aprovando provas
        gráfica e conteúdo;
  \item
        fazer diagramação de livros, jornais e publicações similares;
  \item
        realizar outras atividades correlatadas da unidade de lotação.
\end{enumerate}

Compete às
\href{https://transparencia.cfp.org.br/crp06/wp-content/uploads/sites/29/2023/03/Resolucao-CRP-002-2023.pdf}{\textbf{profissionais
    de suporte administrativo}}
\href{https://transparencia.cfp.org.br/crp06/wp-content/uploads/sites/29/2022/09/Resolucao-03-22-PECS-Anexo-II-Regulamentacao-Plano-de-Empregos-Carreiras-e-Salarios-f-assinada.pdf}{\textbf{(assistente
    administrativo)}}:

\begin{enumerate}

  \item
        executar trabalhos de recepção e suporte administrativo que envolvam:
        serviços de informação e controle de visitantes e correspondências,
        redação de documentos e relatórios, digitação, reprodução xerográfica,
        coleta, expedição, distribuição e arquivamento de documentos;
  \item
        realizar visitas de expedientes junto a fóruns e comarcas;
  \item
        requisitar, receber e controlar o uso de material de consumo;
  \item
        realizar atividades de atendimento ao profissional, bem como
        atividades de caráter administrativo que envolva redação e digitação;
  \item
        preparar e dar suporte administrativo às sessões plenárias éticas e de
        julgamento;
  \item
        atender às solicitações das comissões técnicas ou grupos de trabalho
        específicos criados pelo CRP-06;
  \item
        efetuar o registro das solicitações, interpretar e encaminhar ao
        demandante, como subsídios para as ações do CRP-06;
  \item
        manter sistema de informação atualizado de todas as solicitações
        recebidas e encaminhadas para conhecimento interno e da direção do
        CRP-06;
  \item
        verificar a situação profissional dos psicólogos contratados por
        instituições ou empresas que prestem serviços de psicologia;
  \item
        fazer atendimento presencial e por telefone aos usuários dos serviços
        do CRP~SP;
  \item
        preparar, organizar, atualizar, classificar, arquivar, elaborar,
        redigir, conferir e expedir documentos que envolva processos de
        registro, transferência, cancelamento da categoria de Psicólogos;
  \item
        informar processos e fornecer subsídios para análise e tomada de
        decisão da diretoria;
  \item
        controlar e manter atualizado o arquivamento dos prontuários de
        Psicólogos;
  \item
        dar suporte administrativo às seções plenárias ordinárias, diretoria,
        comissões temáticas, grupos de trabalho e núcleos;
  \item
        providenciar o deslocamento e hospedagem dos conselheiros de plenário,
        diretoria, grupos de trabalho, comissões e núcleos;
  \item
        preparar, organizar, atualizar, elaborar e conferir documentos,
        relatórios, minutas de atos administrativos e normativos, gráficos e
        tabelas;
  \item
        executar rotinas, sistemas, métodos e procedimentos de trabalho
        relativos ao processamento de folha de pagamentos, recolhimento de
        encargos sociais, rotinas de administração de pessoal, recrutamento e
        seleção e administração de benefícios;
  \item
        executar processos de admissão, rescisão e administração de contratos
        de trabalho;
  \item
        executar rotinas, sistemas, métodos e procedimentos de trabalho
        relativos a administração de fornecedores, contratos públicos,
        materiais, processo de compras, recebimento e expedição de materiais,
        controle de estoque, análise de custos, preços e resultados, análise
        de viabilidade econômica, administração e controle orçamentário;
  \item
        preparar, organizar, atualizar, elaborar e conferir documentos,
        relatórios, gráficos e tabelas bem como prestar atendimento ao
        profissional e ao público em geral;
  \item
        controlar fundo fixo (pequeno caixa), verbas, contas a pagar e fluxo
        de caixa, conferindo notas fiscais/recibos e prestando contas;
  \item
        realizar conciliação e análise de movimentos contábeis;
  \item
        fazer lançamentos contábeis da folha de pagamentos, balancetes e
        outros demonstrativos de contas preparando / elaborando a declaração
        de imposto de renda, DIRF e demonstrativos gerenciais, cadastramento e
        controle do arquivo inativo e do inventário físico;
  \item
        realizar outras atividades correlatas da unidade de lotação.
\end{enumerate}